% ============================================================
%  06-introduction.tex — Introduction Générale
% ============================================================
\chapterstar{Introduction Générale}

\vspace{-6pt}
\begin{tcolorbox}[
  enhanced, colback=sectionbg, colframe=isetblue!40,
  arc=4pt, boxrule=0pt,
  borderline west={4pt}{0pt}{isetgold},
  left=10pt, right=10pt, top=6pt, bottom=6pt,
  fontupper=\small\itshape\color{darkgray}
]
Ce rapport présente le travail réalisé dans le cadre du Projet de Fin d'Études
(PFE) pour l'obtention de la Licence en \textbf{Développement des Systèmes
d'Information} à l'ISET Nabeul, année universitaire~2024/2025.
Il documente la conception et la réalisation complète de la plateforme
\textbf{SolarEase}, depuis l'analyse des besoins jusqu'au déploiement.
\end{tcolorbox}

\vspace{12pt}

% ──────────────────────────────────────────────────────────
\section*{Contexte et motivation}

La transition énergétique figure parmi les enjeux stratégiques majeurs du
XXI\textsuperscript{e} siècle. En Tunisie, la filière photovoltaïque connaît
une croissance soutenue, portée par deux dynamiques convergentes~:
d'une part, les objectifs nationaux d'efficacité énergétique inscrits dans
le \textit{Plan Solaire Tunisien} (30\,\% d'énergies renouvelables à
l'horizon 2030)~; d'autre part, la baisse continue des coûts des modules
solaires qui rend cette technologie accessible à un public de plus en plus
large --- résidentiel, commercial et industriel.

Des centaines de sociétés d'installation photovoltaïque opèrent aujourd'hui
sur le territoire tunisien. Ces acteurs jouent un rôle clé dans la
concrétisation de la transition énergétique, mais peinent encore à
industrialiser leurs processus internes~: dimensionnement, suivi client,
facturation et génération de rapports techniques.

% ──────────────────────────────────────────────────────────
\section*{Problématique}

Une analyse de terrain révèle que la majorité de ces installateurs s'appuient
sur des outils hétérogènes et non intégrés~: tableurs manuels pour le
dimensionnement, logiciels isolés pour la simulation, courriers électroniques
pour la relation client et documents bureautiques pour les devis. Cette
fragmentation engendre~:
\begin{itemize}
  \item des erreurs de calcul lors du dimensionnement~;
  \item une perte de temps considérable sur les tâches répétitives~;
  \item l'absence de traçabilité des projets~;
  \item une image peu professionnelle vis-à-vis du client final.
\end{itemize}

\begin{tcolorbox}[
  enhanced, colback=isetblue!7, colframe=isetblue,
  arc=4pt, boxrule=0.6pt,
  left=12pt, right=12pt, top=10pt, bottom=10pt,
  fontupper=\normalsize\itshape\color{isetblue!80!black}
]
\textbf{Problématique centrale~:} Comment concevoir et réaliser une plateforme
web unifiée permettant aux installateurs photovoltaïques de gérer l'intégralité
de leur cycle métier --- du premier contact client jusqu'au rapport de
dimensionnement --- en s'appuyant sur des données fiables (PVGIS) et des
recommandations intelligentes (IA), tout en garantissant sécurité, scalabilité
et évolutivité~?
\end{tcolorbox}

% ──────────────────────────────────────────────────────────
\section*{Solution proposée~: SolarEase}

Pour répondre à cette problématique, nous proposons \textbf{SolarEase}~---~une
plateforme web intelligente dédiée aux installateurs photovoltaïques. Elle
centralise l'intégralité du cycle de vie d'un projet solaire au sein d'une
interface unique, depuis la gestion du client jusqu'à la remise d'un rapport
PDF professionnel, en passant par le calcul de dimensionnement et l'analyse
financière sur 25~ans.

SolarEase se distingue par trois innovations majeures~:
\begin{enumerate}
  \item \textbf{Données réelles PVGIS}~: les calculs de production solaire
    s'appuient sur l'API officielle de la Commission Européenne (données
    d'irradiance géolocalisées), garantissant une précision scientifique~;
  \item \textbf{Mode Night~Panel}~: une fonctionnalité inédite permettant
    de simuler et comparer des installations avec stockage batterie intégré,
    incluant le calcul de la couverture nocturne et du taux
    d'autoconsommation~;
  \item \textbf{Intelligence Artificielle locale}~: des recommandations
    personnalisées générées par un LLM (Ollama / TinyLLaMA) déployé
    localement, sans dépendance à des services cloud externes.
\end{enumerate}

% ──────────────────────────────────────────────────────────
\section*{Méthodologie et organisation du rapport}

Le projet a été mené selon la méthodologie \textbf{Kanban}, offrant un flux
continu de développement piloté visuellement via \textbf{Jira}. Le travail
est découpé en \textbf{cinq sprints} itératifs de deux semaines chacun,
permettant des livraisons incrémentales et une validation régulière avec
l'encadrant.

Le présent rapport est structuré comme suit~:

\begin{table}[H]
  \caption*{Structure du rapport}
  \centering\small
  \setlength{\tabcolsep}{8pt}
  \begin{tabularx}{\textwidth}{L{4.5cm} X}
    \toprule
    \textbf{Section} & \textbf{Contenu} \\
    \midrule
    Chapitre 1 --- Cadre général &
    Présentation de l'organisme d'accueil, problématique, objectifs,
    analyse critique de l'existant et choix méthodologique \\
    \addlinespace
    Chapitre 2 --- Besoins \& Architecture &
    Spécification des besoins, pilotage Kanban, Product Backlog,
    planification des sprints et architecture technique globale \\
    \addlinespace
    Chapitre 3 --- Sprints 1 et 2 &
    Fondation technique (authentification JWT/OTP, Gateway) et gestion
    des clients et des projets (CRUD, cycle de vie, tableau de bord) \\
    \addlinespace
    Chapitre 4 --- Sprints 3, 4 et 5 &
    Dimensionnement solaire (PVGIS, Night Panel), approche RAG \textbf{hybride}
    (catalogue société + \texttt{pgvector}) pour les recommandations IA,
    génération PDF, suivi terrain, messagerie
    temps réel, notifications, espace client, demandes, devis,
    factures et workflow commercial complet \\
    \addlinespace
    Chapitre 5 --- DevOps \& mise en production &
    Containerisation des microservices backend, pipeline CI/CD GitHub
    Actions, observabilité (Prometheus, Grafana, Loki), qualité et
    sécurité (JaCoCo, Trivy, Gitleaks), Infrastructure as Code
    (Kubernetes + Kustomize) et stratégie de déploiement
    multi-environnements (Railway pour le backend, Vercel pour le frontend) \\
    \addlinespace[8pt]
    Conclusion générale &
    Bilan des réalisations, limites et perspectives d'évolution \\
    \bottomrule
  \end{tabularx}
\end{table}

\cleardoublepage
